% ============================================================
%  Informe Ejercicios Opcionales Práctica 2 — SOII
%  Compilar: lualatex informe_opcionales.tex  (dos pasadas)
% ============================================================
\documentclass[a4paper, 12pt]{article}

% ── Codificación y fuentes ──────────────────────────────────
\usepackage[T1]{fontenc}
\usepackage[utf8]{inputenc}

% ── Idioma ──────────────────────────────────────────────────
\usepackage[spanish, es-nolists, es-noindentfirst]{babel}

% ── Geometría ───────────────────────────────────────────────
\usepackage[
  a4paper,
  top=2.5cm, bottom=2.5cm,
  left=2.8cm, right=2.8cm
]{geometry}

% ── Colores ─────────────────────────────────────────────────
\usepackage{xcolor}
\definecolor{azul}     {HTML}{2E5FA3}
\definecolor{azulosc}  {HTML}{1A3A6B}
\definecolor{grisclaro}{HTML}{F4F4F4}
\definecolor{grisborde}{HTML}{CCCCCC}
\definecolor{rojocod}  {HTML}{A31515}

% ── Encabezados y pies ──────────────────────────────────────
\usepackage{fancyhdr}
\pagestyle{fancy}
\fancyhf{}
\fancyhead[L]{\small\color{azul}\textbf{SOII --- Práctica 2}}
\fancyhead[R]{\small\color{gray}Ejercicios Opcionales (Hilos)}
\fancyfoot[C]{\small\thepage}
\renewcommand{\headrulewidth}{0.4pt}

% ── Títulos de sección ──────────────────────────────────────
\usepackage{titlesec}
\titleformat{\section}
  {\Large\bfseries\color{azul}}
  {\thesection.}{0.6em}{}
  [\vspace{-0.3em}\textcolor{azul}{\rule{\linewidth}{0.8pt}}\vspace{0.2em}]

\titleformat{\subsection}
  {\large\bfseries\color{azulosc}}
  {\thesubsection.}{0.5em}{}

\titleformat{\subsubsection}
  {\normalsize\bfseries}
  {\thesubsubsection.}{0.4em}{}

% ── Listas ──────────────────────────────────────────────────
\usepackage{enumitem}
\setlist[itemize]{leftmargin=1.5em, itemsep=2pt, topsep=4pt}

% ── Tablas ──────────────────────────────────────────────────
\usepackage{booktabs}
\usepackage{array}
\usepackage{colortbl}
\newcolumntype{L}[1]{>{\raggedright\arraybackslash}p{#1}}
\newcolumntype{C}[1]{>{\centering\arraybackslash}p{#1}}

% ── Código fuente ───────────────────────────────────────────
\usepackage{listings}
\lstdefinestyle{cstyle}{
  language=C,
  basicstyle=\ttfamily\small,
  keywordstyle=\bfseries\color{azul},
  commentstyle=\itshape\color{gray},
  stringstyle=\color{rojocod},
  numberstyle=\tiny\color{gray},
  numbers=left,
  numbersep=8pt,
  stepnumber=1,
  backgroundcolor=\color{grisclaro},
  frame=single,
  framesep=4pt,
  rulecolor=\color{grisborde},
  breaklines=true,
  tabsize=4,
  showstringspaces=false,
  xleftmargin=1.8em,
  xrightmargin=0.5em,
  captionpos=b,
  literate=
    {á}{{\'a}}1 {é}{{\'e}}1 {í}{{\'i}}1 {ó}{{\'o}}1 {ú}{{\'u}}1
    {Á}{{\'A}}1 {É}{{\'E}}1 {Í}{{\'I}}1 {Ó}{{\'O}}1 {Ú}{{\'U}}1
    {ñ}{{\~n}}1 {Ñ}{{\~N}}1,
}
\lstdefinestyle{shellstyle}{
  basicstyle=\ttfamily\small,
  backgroundcolor=\color{grisclaro},
  frame=single,
  framesep=4pt,
  rulecolor=\color{grisborde},
  breaklines=true,
  numbers=none,
  xleftmargin=1em,
  xrightmargin=0.5em,
  literate={→}{{$\rightarrow$}}1 {}{{\~n}}1,
}
\lstset{style=cstyle}

% ── Cajas resaltadas ────────────────────────────────────────
\usepackage{tcolorbox}
\tcbuselibrary{skins, breakable}

\newtcolorbox{infobox}[1][]{
  colback=azul!7, colframe=azul, fonttitle=\bfseries,
  title=#1, left=6pt, right=6pt, top=4pt, bottom=4pt,
  breakable
}
\newtcolorbox{warnbox}[1][]{
  colback=orange!8, colframe=orange!70!black, fonttitle=\bfseries,
  title=#1, left=6pt, right=6pt, top=4pt, bottom=4pt,
  breakable
}

% ── Hiperenlaces ────────────────────────────────────────────
\usepackage{hyperref}
\hypersetup{
  colorlinks=true, linkcolor=azul, urlcolor=azul,
  pdftitle={Práctica 2 SOII --- Ejercicios Opcionales},
}

% ── Miscelánea ──────────────────────────────────────────────
\usepackage{microtype}
\usepackage{parskip}
\setlength{\parskip}{6pt}

% ── Macro inline-code ───────────────────────────────────────
\newcommand{\ic}[1]{\texttt{#1}}

% ============================================================
\begin{document}

% ── Portada ─────────────────────────────────────────────────
\begin{titlepage}
  \centering
  \vspace*{3cm}
  {\Huge\bfseries\color{azul} Sistemas Operativos II\par}
  \vspace{0.6cm}
  {\LARGE Práctica 2 --- Informe Opcionales\par}
  \vspace{0.4cm}
  {\large\color{gray} Sincronización avanzada con hilos POSIX (pthreads)\par}
  \vspace{0.5cm}
  \textcolor{azul}{\rule{0.7\linewidth}{1.2pt}}
  \vspace{2cm}

  \begin{tabular}{rl}
    \textbf{Asignatura:} & Sistemas Operativos II \\[4pt]
    \textbf{Titulación:} & Grado en Ingeniería Informática \\[4pt]
    \textbf{Curso:}      & 2024/25 \\[4pt]
    \textbf{Fecha:}      & \today \\
  \end{tabular}

  \vfill
  {\small Universidad de Santiago de Compostela}
\end{titlepage}

% ── Índice en la segunda página ─────────────────────────────
\tableofcontents
\newpage

% ============================================================
\section{Ejercicio Opcional 1: Productor-Consumidor con Hilos}
% ============================================================

\subsection{Descripción arquitectónica}

A diferencia de los apartados obligatorios donde se utilizaban procesos pesados comunicados mediante \ic{mmap()}, en esta versión (\ic{opc1\_threads.c}) se instancian dos hilos dentro de un único proceso utilizando la biblioteca POSIX Threads (\ic{pthread}).

La memoria es compartida de forma natural ya que todos los hilos residen en el mismo espacio de direcciones virtuales. Esto elimina la necesidad de utilizar ficheros mapeados en memoria.

\begin{infobox}[Semáforos sin nombre (\texttt{sem\_init})]
  En lugar de crear semáforos nombrados en el sistema de archivos del kernel mediante \ic{sem\_open()}, se utilizan semáforos anónimos basados en memoria mediante \ic{sem\_init(\&sem, 0, valor)}. El segundo argumento a \ic{0} indica que el semáforo será compartido exclusivamente entre los hilos del proceso actual, evitando dejar recursos huérfanos en el sistema.
\end{infobox}

\subsection{Análisis de resultados}

La ejecución evidencia tres fases temporales claramente delimitadas por los \ic{sleep()} situados en el código:
\begin{enumerate}
    \item \textbf{Productor más rápido:} El buffer se llena rápidamente hasta su capacidad máxima (\ic{top=10}).
    \item \textbf{Consumidor más rápido:} El consumidor agota los recursos del buffer llegando a \ic{top=0}.
    \item \textbf{Fase aleatoria:} Ambos procesos se bloquean alternativamente.
\end{enumerate}

La salida final corrobora el correcto uso del semáforo mutex para la sección crítica:

\begin{lstlisting}[style=shellstyle]
[MAIN] === Vocales producidas ===
  a/A=8  e/E=6  i/I=8  o/O=5  u/U=3
[MAIN] === Vocales consumidas ===
  a/A=8  e/E=6  i/I=8  o/O=5  u/U=3
[MAIN] OK: vocales producidas == vocales consumidas.
\end{lstlisting}

\newpage

% ============================================================
\section{Ejercicio Opcional 2: $N$ Productores y $M$ Consumidores}
% ============================================================



\subsection{Diseño y control de concurrencia}

La generalización del problema a múltiples productores y consumidores (\ic{opc2\_multi.c}) introduce complicaciones de sincronización adicionales. Además de proteger el buffer, es necesario proteger otros recursos compartidos:

\begin{itemize}
    \item \textbf{\ic{mutex\_prod}:} Protege el acceso al puntero de lectura del fichero \ic{input.txt} compartido entre los $N$ productores.
    \item \textbf{\ic{mutex\_fin}:} (Junto con variables \ic{atomic\_int}) Protege los contadores globales de elementos producidos y consumidos para gestionar la terminación.
    \item \textbf{\ic{mutex\_vocales\_[prod|cons]}:} Evita condiciones de carrera al actualizar los contadores estadísticos de vocales.
\end{itemize}

\subsection{Protocolo de Terminación en Redes M:N}

El problema clásico en redes de varios productores y consumidores es evitar que los consumidores queden bloqueados en un \ic{sem\_wait(\&llenas)} infinito una vez que los productores han terminado.

La solución implementada utiliza una combinación de variables atómicas y \emph{broadcasting} mediante semáforos:
\begin{enumerate}
    \item El hilo principal (\ic{main}) espera mediante \ic{pthread\_join()} a que todos los productores terminen.
    \item Se activa el flag atómico \ic{produccion\_fin}.
    \item El hilo principal ejecuta \ic{sem\_post(\&llenas)} $M$ veces (una por cada consumidor).
    \item Los consumidores atrapados en espera se despiertan, evalúan la condición de finalización y propagan la señal si es necesario antes de hacer \ic{break}.
\end{enumerate}

\subsection{Análisis de resultados}

La ejecución solicitada de 3 productores y 2 consumidores reporta una sincronización perfecta para 120 elementos en total (40 elementos por productor):

\begin{lstlisting}[style=shellstyle]
[MAIN] 3 productor(es), 2 consumidor(es)
[MAIN] Total a producir/consumir: 120 items
...
[PROD-0] i=39  producido='c'
[PROD-0] i=39  insertado  top=5  producidos_total=120
[CONS-0] retirado='c'  top=4  consumidos_total=116
...
[MAIN] === Resumen ===
[MAIN] Producidos : 120
[MAIN] Consumidos : 120
[MAIN] Vocales producidas: a=12 e=12 i=8 o=8 u=4
[MAIN] Vocales consumidas: a=12 e=12 i=8 o=8 u=4
[MAIN] OK: producidos == consumidos, vocales coinciden.
\end{lstlisting}

\newpage

% ============================================================
\section{Ejercicio Opcional 3: Arquitectura en Pipeline}
% ============================================================



\subsection{Descripción del Pipeline Lineal}

El archivo \ic{opc3\_pipeline.c} implementa un flujo de datos secuencial a través de $P$ hilos, conectados por $P-1$ buffers intermedios de tamaño fijo. 
Cada hilo intermedio actúa simultáneamente como \textbf{consumidor} de la etapa anterior y como \textbf{productor} de la etapa siguiente. 

La topología permite el \textbf{máximo grado de concurrencia}: al disponer de arrays independientes de semáforos (\ic{g\_vacias[k]}, \ic{g\_llenas[k]}, \ic{g\_mutex[k]}), el hilo 1 puede estar escribiendo en el buffer 1 al mismo tiempo que el hilo 2 escribe en el buffer 2. No existe un \emph{Global Interpreter Lock} ni bloqueo general del sistema.

\subsection{Transformación de datos (\texttt{next\_alpha})}

Para evidenciar el flujo a través de las etapas intermedias, cada hilo aplica una función de cifrado de rotación unitaria cíclica (tipo César $+1$) exclusivamente a los caracteres alfabéticos.

\subsection{Propagación en cascada y resultados}

Para terminar la ejecución, el Productor Origen (Hilo 0) inserta un carácter centinela nulo (\ic{'\textbackslash 0'}) que cada hilo intermedio detecta y propaga a la etapa adyacente antes de terminar su propia ejecución.

Se muestra un pipeline configurado con $P=4$ etapas (2 etapas intermedias de transformación):

\begin{lstlisting}[style=shellstyle]
[MAIN] Pipeline de 4 etapas, 3 buffer(es) intermedios.
[MAIN] Cada caracter pasa por 2 transformacion(es) next_alpha.
[MAIN] 80 items a producir.
...
[Hilo  0] i= 0  leido='E'
[Hilo  0->buf 0] insertado='E'  top=1
[buf 0->Hilo  1] retirado='E'  top=0
[Hilo  1] 'E' -> 'F'
[Hilo  1->buf 1] insertado='F'  top=1
[buf 1->Hilo  2] retirado='F'  top=0
[Hilo  2] 'F' -> 'G'
[Hilo  2->buf 2] insertado='G'  top=1
[buf 2->Hilo  3] retirado='G'  top=0
[Hilo  3] consumido='G'
...
[MAIN] === Resumen del pipeline ===
[MAIN] Transformaciones aplicadas: 2 (next_alpha x2)
[MAIN] Vocales en origen  (hilo 0):   a=8 e=6 i=8 o=5 u=3
[MAIN] Vocales al final   (hilo 3):  a=1 e=3 i=2 o=2 u=0
[MAIN] (Las vocales pueden cambiar tras transformaciones: correcto)
[MAIN] Vocales totales origen=30, finales=8
[MAIN] OK: se procesaron exactamente 80 caracteres.
\end{lstlisting}

\begin{warnbox}[Discrepancia de Vocales]
  A diferencia de los ejercicios anteriores, el conteo de vocales inicial y final no coincide. Esto es estrictamente correcto y esperado, ya que las transformaciones intermedias (ej. $e \rightarrow f \rightarrow g$) convierten vocales en consonantes y viceversa. La validación de integridad recae en que los caracteres iniciales y procesados totales son exactamente 80.
\end{warnbox}

\end{document}